%%==================================================================
\section{Introduction}\label{sec:intro}
%%==================================================================

\subsection{The theorem of \oai{}}
On 25 October 2026 \oai{} released a preprint \cite{OAI26} with a deterministic
algorithm that, for every fixed $k\ge 2$, counts the satisfying assignments of a
$k$-CNF formula with $n$ variables and $m$ clauses in time
\begin{equation}\label{eq:oaitime}
  (n+m)^{O_k(1)}\, 2^{O_k(n\log\log n/\sqrt{\log n})}
\end{equation}
and polynomial space. Since $2^{O(n\log\log n/\sqrt{\log n})}=2^{o(n)}$, the
exponential time hypothesis of Impagliazzo and Paturi \cite{IP01}, which asserted
that $3$-SAT needs time $2^{\Omega(n)}$, is false, and with it the strong
hypothesis \cite{IPZ01}, the counting hypothesis $\#$ETH \cite{DHMTW14}, and the
gap and nondeterministic variants. The algorithm is algebraic. It writes the
count as a sum over the Boolean cube of a product of local polynomials, replaces
every variable occurrence by a linear form in a parameter $T$ chosen so that at
each of $Q$ sample points the factors separate into $r$ groups after a short
list of integer measurements of the assignment has been fixed, recovers the
Taylor coefficients of the resulting univariate polynomial at the sample points
from Boolean substitutions, and interpolates back to $T=0$. The present paper
owes its existence to that preprint, and every step of the algorithm that is used
below is \oai{}'s.\footnote{The preprint is unrefereed. For this reason the
algorithm and its analysis are reproduced here in full, in Sections
\ref{sec:routine} and \ref{sec:design}, so that no statement of this paper
depends on a text outside it. The reproduction follows the preprint closely, with
the notation of this paper, and the reader will find the corresponding statements
of \cite{OAI26} named in footnotes. The preprint also contains a Lean~4
formalisation of its main theorem; it is not used here.}

\subsection{What this paper adds}
The counting routine of \cite{OAI26} is stated there for factors with rational
coefficients, with the remark that it works over any commutative $\Q$-algebra.
We take that remark seriously and run the algorithm over truncated polynomial
rings $\Q[W]/(W^{N})$. A factor may then be $W^{\,w(x)}$ for an integer weight
function $w$ of a few variables, and the sum over the cube becomes a generating
polynomial: its coefficient of $W^{j}$ counts the assignments of total weight
$j$. One number has turned into a profile, and the profile contains an optimum.
The cost of the change is nil in the exponent, and polynomial in $N$ elsewhere.

The technical statement is Theorem~\ref{thm:weighted}: for fixed $k$ and $c$, a
system of at most $cn$ factors, each depending on at most $k$ of the $n$ Boolean
variables and taking values in a commutative $\Q$-algebra $\A$, is summed over
$\{0,1\}^n$ in
\begin{equation}\label{eq:rho}
  2^{O_{k,c}(\rhon(n))}\poly(n) \quad\text{operations in $\A$},\qquad
  \rhon(n)=\frac{n(\log\log n)^2}{\sqrt{\log n}} .
\end{equation}
The exponent is the one \oai{} obtain in their Sections 3--6 before the
refinement that removes one logarithmic factor; we do not reproduce that
refinement, because none of the consequences below needs it.

From Theorem~\ref{thm:weighted} we deduce, with complete proofs:
\begin{enumerate}[label=(\roman*),leftmargin=2.2em]
\item\label{it:profile} for every fixed $k$ and $c$, the \emph{profile} of a
constraint satisfaction instance with $n$ Boolean variables and at most $cn$
constraints of arity at most $k$, that is the number of assignments satisfying
exactly $j$ constraints for every $j$, in time $2^{O_{k,c}(\rhon(n))}$
(Theorem~\ref{thm:profile}); in particular the exact optimum of Max-$k$-SAT and
of Max-Cut on sparse instances, together with the number of optimal solutions
(Corollaries \ref{cor:maxsat} and \ref{cor:maxcut});
\item\label{it:indep} the independence polynomial of \emph{every} graph on $n$
vertices, hence its independence number, clique number, vertex cover number and the
number of maximum independent sets, in time
$2^{O(n(\log\log n)^2/(\log n)^{1/4})}$ (Theorem~\ref{thm:indep}); the proof
combines Theorem~\ref{thm:weighted} with a branching rule on vertices of large
degree whose cost we bound explicitly;
\item\label{it:vc} Vertex Cover with parameter $k$ in time
$2^{O(k(\log\log k)^2/(\log k)^{1/4})}+\poly(n)$, that is $2^{o(k)}$ up to a
polynomial, through the kernel of Nemhauser and Trotter (Theorem~\ref{thm:vc});
the lower bound $2^{\Omega(k)}$, which was a standard consequence of the
exponential time hypothesis \cite[Theorem 14.6]{CFKLMPPS15}, therefore fails;
\item\label{it:ising} for every fixed $\Delta$, the density of states of the
Ising model with integer couplings and fields, and of the $q$-state Potts model,
on graphs of maximum degree $\Delta$ with $n$ vertices, in time
$2^{O_{\Delta,q}(\rhon(n))}$ (Theorem~\ref{thm:ising}); every partition function
and every ground state follows;
\item\label{it:bounded} the minimum dominating set and the matching polynomial of
a bounded-degree graph, the permanent of a sparse integer matrix, and the
minimum hitting set and the maximum split of a sparse set system with sets of
bounded size, in time $2^{O(\rhon(n))}$ (Theorems \ref{thm:bounded} and
\ref{thm:permanent});
\item\label{it:domain} for fixed domain size and arity, every sparse constraint
system over a finite domain, hence every partition function $Z_A(G)$ of a
fixed symmetric rational matrix $A$ and the number of homomorphisms into a
fixed graph, on bounded-degree graphs, in time $2^{O(\rhon(n))}$
(Theorem~\ref{thm:domain}, Corollary~\ref{cor:hom});
\item\label{it:global} the satisfying assignments of a sparse $k$-CNF formula
of each Hamming weight, and those whose weight lies in a prescribed residue
class, in time $2^{O(\rhon(n))}$ (Corollary~\ref{cor:hamming},
Proposition~\ref{prop:modular});
\item\label{it:snp} $k$-SAT lies in the class $\SE$ of problems solvable in
time $2^{\varepsilon n}\poly$ for every $\varepsilon>0$, by the sparsification
lemma of Impagliazzo, Paturi and Zane \cite{IPZ01}, and therefore every problem of
the class $\SNP$ of Papadimitriou and Yannakakis \cite{PY91} lies in $\SE$
(Theorem~\ref{thm:snp}).
\end{enumerate}
Items \ref{it:profile}--\ref{it:global} concern problems for which no
$2^{o(n)}$ algorithm was known and for which such an algorithm was excluded under
the exponential time hypothesis by standard reductions. The best exact algorithms
for the maximum independent set run in time $2^{0.1996n}$ \cite{XN17}, for
Max-Cut in time $2^{\omega n/3}$ where $\omega$ is the exponent of matrix
multiplication \cite{Wil05}, and for the permanent and the perfect matchings of
sparse graphs in time $2^{n-\Omega(n/\Delta)}$ \cite{CP15}; each of these bounds
is exceeded here asymptotically.

\subsection{What this paper does not claim}
The algorithm counts sums of products of \emph{local} factors over a cube of
dimension $n$ and its running time is subexponential in $n$. Two features limit
its reach, and they are stated here so that no reader expects more.
First, the exponent is $o(n)$ and not $n^{o(1)}$. A polynomial-time reduction
from an NP-complete problem to $3$-SAT multiplies the number of variables by a
polynomial, so the algorithm says nothing about $\mathrm{P}$ versus
$\mathrm{NP}$, and nothing here bears on that question. Second, the number of
factors and the total degree of their product must be $O(n)$, after a branching
step that fixes $O(n/\log n)$ variables. Dense Max-Cut, the chromatic polynomial,
the Tutte polynomial, Orthogonal Vectors, $3$SUM and All-Pairs Shortest Paths
need either $\Theta(n^2)$ constraints or reductions that multiply the variable
count, and they are untouched. The paper contains no partial result on any of
these problems.

\subsection{Organisation}
Section~\ref{sec:prelim} fixes notation and the model of computation over an
algebra. Sections \ref{sec:routine} and \ref{sec:design} reproduce the
summation routine and the construction of the replacement forms. Section
\ref{sec:weighted} proves Theorem~\ref{thm:weighted}, including the branching on
frequent variables and the bit-complexity bound over $\Q[W]/(W^N)$. Section
\ref{sec:applications} proves items \ref{it:profile}--\ref{it:global}, Section
\ref{sec:snp} proves item \ref{it:snp}, and Section~\ref{sec:scope} records what
the method does not give. Appendix~\ref{app:check} describes an exact-arithmetic
check of the extended algorithm on small instances; no statement of the paper
depends on it. The author's earlier work on algorithmic structures in geometric and
topological settings \cite{Bha26a,Bha26b,Bha26c} is unrelated in method but
shares the aim of making a computation explicit enough to be verified line by
line.

A word on the constants. The parameters of Lemma~\ref{lem:params} become
admissible only when $\sqrt{\log n}$ exceeds a large multiple of
$(\log\log n)^2$, that is for $n$ beyond any practical size; the same is true
of \oai{}'s algorithm, and Figure~\ref{fig:costs} shows where the three costs
of the routine settle. The theorems are statements about asymptotic complexity,
and the paper does not present them as anything else.
